%%%%%%%%%%%%%%%%%%%%%%%%%%%%%%%%%%%%%%%%%
% Academic Title Page
% LaTeX Template
% Version 2.0 (17/7/17)
%
% This template was downloaded from:
% http://www.LaTeXTemplates.com
%
% Original author:
% WikiBooks (LaTeX - Title Creation) with modifications by:
% Vel (vel@latextemplates.com)
%
% License:
% CC BY-NC-SA 3.0 (http://creativecommons.org/licenses/by-nc-sa/3.0/)
% 
% Instructions for using this template:
% This title page is capable of being compiled as is. This is not useful for 
% including it in another document. To do this, you have two options: 
%
% 1) Copy/paste everything between \begin{document} and \end{document} 
% starting at \begin{titlepage} and paste this into another LaTeX file where you 
% want your title page.
% OR
% 2) Remove everything outside the \begin{titlepage} and \end{titlepage}, rename
% this file and move it to the same directory as the LaTeX file you wish to add it to. 
% Then add \input{./<new filename>.tex} to your LaTeX file where you want your
% title page.
%
%%%%%%%%%%%%%%%%%%%%%%%%%%%%%%%%%%%%%%%%%

%----------------------------------------------------------------------------------------
%	PACKAGES AND OTHER DOCUMENT CONFIGURATIONS
%----------------------------------------------------------------------------------------

\documentclass[12pt]{article}

\usepackage[utf8]{inputenc} % Required for inputting international characters
\usepackage[T1]{fontenc} % Output font encoding for international characters

\usepackage{mathpazo} % Palatino font

\usepackage{comment} % enables the use of multi-line comments (\ifx \fi) 
\usepackage{lipsum} %This package just generates Lorem Ipsum filler text. 
% \usepackage{fullpage} % changes the margin
\usepackage{CJKutf8}
\usepackage{enumitem}
\usepackage{titlesec}
\usepackage[english]{babel}
\usepackage{blindtext}
\usepackage{graphicx}     % for figure
\usepackage{subcaption}   % for figure
\usepackage[export]{adjustbox}
\usepackage[most]{tcolorbox}
\usepackage{xcolor}
\usepackage{geometry}

\newenvironment{reportsection}[1]{%
\section*{#1}%
\vspace{-1em}
\begin{enumerate}[label=\alph*.,labelindent=\parindent, leftmargin=0em]{%
% \setlength{\topmargin}{0cm}%
% \setlength{\topsep}{0pt}%
% \setlength{\leftmargin}{0cm}%
% \setlength{\rightmargin}{0cm}%
% \setlength{\listparindent}{\parindent}%
% \setlength{\itemindent}{\parindent}%
% \setlength{\parsep}{\parskip}%
}%
\item[]}{\end{enumerate}}

\titlespacing*{\section} {-1.5em}{2.5ex plus 1ex minus .2ex}{1.3ex plus .2ex}
\titleformat{\section}
{\normalfont\fontsize{18}{15}\bfseries}{\thesection}{1em}{}

\titlespacing*{\subsection}
{0pt}{0.1\baselineskip}{0.1\baselineskip}

\begin{document}

%----------------------------------------------------------------------------------------
%	TITLE PAGE
%----------------------------------------------------------------------------------------

\begin{titlepage} % Suppresses displaying the page number on the title page and the subsequent page counts as page 1
	\newcommand{\HRule}{\rule{\linewidth}{0.5mm}} % Defines a new command for horizontal lines, change thickness here
		
	\center % Centre everything on the page
		
	%------------------------------------------------
	%	Headings
	%------------------------------------------------
		
	\textsc{\LARGE National Taipei University of Technology}\\[1.5cm] % Main heading such as the name of your university/college
		
	\textsc{\Large 2018 Fall}\\[0.5cm] % Major heading such as course name
		
	\textsc{\large 245765 - Advanced Digital Image Processing}\\[0.5cm] % Minor heading such as course title
		
	%------------------------------------------------
	%	Title
	%------------------------------------------------
		
	\HRule\\[0.4cm]
		
	{\huge\bfseries HW\#6-2\\\vspace{1em} Geometric Distortion Transformation}\\[0.4cm] % Title of your document
		
	\HRule\\[1.5cm]
		
	%------------------------------------------------
	%	Author(s)
	%------------------------------------------------
	\begin{CJK}{UTF8}{bsmi}
		\begin{minipage}{0.4\textwidth}
			\begin{flushleft}
				\large
				\textit{Author}\\
				106368002 張昌祺\\ 
				\textsc{Chang-Qi Zhang} \\
				justin840727@gmail.com % Your name
			\end{flushleft}
		\end{minipage}
		~
		\begin{minipage}{0.4\textwidth}
			\begin{flushright}
				\large
				\textit{Advisor}\\
				電子所 \\
				高立人 副教授 % Supervisor's name
			\end{flushright}
		\end{minipage}
	\end{CJK}
	% If you don't want a supervisor, uncomment the two lines below and comment the code above
	%{\large\textit{Author}}\\
	%John \textsc{Smith} % Your name
		
	%------------------------------------------------
	%	Date
	%------------------------------------------------
		
	\vfill\vfill\vfill % Position the date 3/4 down the remaining page
		
	{\large\today} % Date, change the \today to a set date if you want to be precise
		
	%------------------------------------------------
	%	Logo
	%------------------------------------------------
		
	\vfill
	\includegraphics[width=0.6\textwidth]{../../logo.jpg}\\[1cm] % Include a department/university logo - this will require the graphicx package
			 
		%----------------------------------------------------------------------------------------
			
		% \vfill % Push the date up 1/4 of the remaining page
			
	\end{titlepage}
	\newgeometry{left=0.8in,right=0.8in,top=1in,bottom=1in}
	\begin{reportsection}{Problem 2~~~Geometric transformation}
		\setcounter{enumi}{1}
		\item Given chessboard\_256.raw and chessboard\_distorted\_256.raw, find 
		the mapping relationship between them. Use this relationship to convert owl\_256.raw 
		to get the result as owl\_transform. Show the output image and discuss your method. 
		(Figure 10\%; Discussion 10\%)
		\subsection*{Ans}
		First, I use Opencv API to detect all corners from both chessboard image.
		\begin{figure}[htbp]
			\centering
			\begin{subfigure}[b!]{0.43\linewidth}
			  \includegraphics[width=\textwidth]{{"../result_img/chessboard_draw_src"}.png}
			  \caption{Chessboard corner detection.}
			\end{subfigure}
			\hspace{1em}
			\begin{subfigure}[b!]{0.43\linewidth}
			  \includegraphics[width=\textwidth]{{"../result_img/chessboard_draw_dst"}.png} 
			  \caption{Chessboard distorted corner detection.}
			\end{subfigure}
			\caption{OpenCV corner detection API result.}
      \label{fig:corner-detection}
		\end{figure}

		Once we get the corners points, we assume the distortion relationship 
		between two set of corners are Equation~\ref{eq:distortion}.
		\begin{equation}
			\begin{split}
				x' &= x (1+k_1r^2+k_2r^4) \\
				y' &= y (1+k_1r^2+k_2r^4)
			\end{split}
			\label{eq:distortion}
		\end{equation}
		
		So that we can easy to get $k_1$ and $k_2$ values by solving linear system
		shows in Equation~\ref{eq:linear-system}. In the homework I use Least 
		squares to solve it. 
		\begin{equation}
		\begin{bmatrix}
				r_{1}^{2}    & r_{1}^{4} \\
				r_{1}^{2}    & r_{1}^{4} \\
				r_{2}^{2}    & r_{2}^{4} \\
				r_{2}^{2}    & r_{2}^{4} \\
				\vdots 			 &\vdots		 \\
				r_{n}^{2}    & r_{n}^{4} \\
				r_{n}^{2}    & r_{n}^{4}
		\end{bmatrix}
		\times
		\begin{bmatrix}
			k_1 \\
			k_2
		\end{bmatrix}
		=
		\begin{bmatrix}
			x_{1}'/x_{1}-1 \\
			y_{1}'/y_{1}-1 \\
			x_{2}'/x_{2}-1 \\
			y_{2}'/y_{2}-1 \\
			\vdots				 \\
			x_{n}'/x_{n}-1 \\
			y_{n}'/y_{n}-1 \\
		\end{bmatrix}
		\label{eq:linear-system}
	\end{equation}
	\newpage
	Once we get $k_1$ and $k_2$ values we can apply to owl\_256.raw image
	to get distorted image. After transformation, it cause some points can
	not map pixel values form original image. Therefore, we need to use 
	nearest pixel value fill into those pixels. The final transformed result
	shows in Figure~\ref{fig:owl-distorted}.
	\begin{figure}[htbp]
		\begin{center}
			\includegraphics[width=0.4\textwidth]{{"../result_img/distored_tf"}.png} 
				\caption{distorted owl\_256.raw result.}
				\label{fig:owl-distorted}
		\end{center} 
	\end{figure}
	\end{reportsection}
	
\end{document}
